\section{Discussion \& Architectural Implications}

The ST-016 results yield three actionable design principles for runtime system
architects. Each principle is grounded in a specific formal or empirical result
from this work.

\begin{enumerate}
    \item \textbf{Safety Contract Decoupling:} Safety verification
    ($C_{\text{contract}}$) cannot be subsumed by statistical policy estimation
    alone. \emph{Grounded in:} EXP-004 ablation of $C_{\text{contract}}$ produced
    EXECUTE on a contract-violating state where the full runtime correctly returned
    STOP. The formal witness artifact was verified consistent with the Lean model.

    \item \textbf{Dynamic Boundary Monitoring:} Operating near ambiguity boundaries
    requires explicit uncertainty estimation ($C_{\text{uncertainty}}$) to prevent
    premature execution. \emph{Grounded in:} EXP-004 ablation of $C_{\text{uncertainty}}$
    at decision margin $M_D \approx 0$ produced EXECUTE where the full runtime
    returned REFINE, with the divergence formally certified via \texttt{RuntimeWitness}.

    \item \textbf{Trajectory History Dependence:} Stateless evaluation fails under
    history-dependent processes, necessitating explicit trajectory prefix tracking
    via $C_{\text{temporal}}$. \emph{Grounded in:} EXP-004 ablation of
    $C_{\text{temporal}}$ produced MONITOR on a history-dependent prefix where
    the full runtime returned INTERVENE (see also Example 1, Section~3).
\end{enumerate}
